% Generated draft: requires mathematical and source review.
\documentclass[11pt]{article}
\usepackage[margin=25mm]{geometry}
\usepackage{fontspec}
\setmainfont{Latin Modern Roman}
\usepackage{amsmath,amssymb,mathtools,unicode-math}
\setmathfont{Latin Modern Math}
\usepackage{enumitem,needspace,xurl,hyperref}
\setlength{\parindent}{0pt}
\setlength{\parskip}{5pt}
\setlength{\emergencystretch}{3em}
\newcommand{\N}{\mathbb{N}}
\newcommand{\Z}{\mathbb{Z}}
\newcommand{\Q}{\mathbb{Q}}
\newcommand{\R}{\mathbb{R}}
\newcommand{\C}{\mathbb{C}}
\newcommand{\F}{\mathbb{F}}
\newcommand{\A}{\mathbb{A}}
\newcommand{\PP}{\mathbb{P}}
\newcommand{\E}{\mathbb{E}}
\newcommand{\eps}{\varepsilon}
\newcommand{\dd}{\,\mathrm{d}}
\begin{document}
\section*{UnsolvedMath ID: 30002218}
\textbf{Source problem code:} catalogue entry
\subsection{Definitions and mathematical statement}
Let \(\mathcal Z\) denote the Jiang--Su \(C^*\)-algebra. A \(C^*\)-algebra \(A\) is separable if it has a countable norm-dense subset. Tensorial absorption of \(\mathcal Z\), also called \(\mathcal Z\)-stability, means that there is a \(C^*\)-algebra isomorphism \(A\cong A\otimes\mathcal Z\).

The supplied question does not define \(d(A,B)\), specify a common ambient algebra, or specify a tensor-product convention. A customary interpretation is the Kadison--Kastler distance for concrete \(C^*\)-subalgebras \(A,B\) of a common \(C^*\)-algebra \(C\). Writing \(A_1=\{a\in A:\|a\|\leq 1\}\) and similarly for \(B_1\), this distance is
\[
 d_{\mathrm{KK}}^C(A,B)=\max\left\{
 \sup_{a\in A_1}\inf_{b\in B_1}\|a-b\|,
 \sup_{b\in B_1}\inf_{a\in A_1}\|b-a\|
 \right\}.
\]
Under this interpretation, and taking \(\otimes\) to mean the spatial \(C^*\)-tensor product, the question is whether
\[
\begin{gathered}
\exists\delta>0\ \ \forall\ C\text{ a }C^*\text{-algebra}\ \ \forall\ A,B\subseteq C\text{ }C^*\text{-subalgebras},\\
\bigl[A\text{ is separable},\quad A\cong A\otimes_{\min}\mathcal Z,\quad d_{\mathrm{KK}}^C(A,B)<\delta\bigr]
\ \Longrightarrow\ B\cong B\otimes_{\min}\mathcal Z.
\end{gathered}
\]
Here \(\delta\) is uniform over the algebras and their common realizations. No separate separability assumption on \(B\), and no unitality or nuclearity assumption on \(A\) or \(B\), is stated in the supplied question. The displayed interpretation of the distance and ambient quantifiers requires confirmation against the original source; it is not a verified transcription of conventions from that report.
\subsection{Short English statement}
The question asks whether one universal positive perturbation bound guarantees that every sufficiently close \(C^*\)-algebra to a separable \(\mathcal Z\)-stable algebra is also \(\mathcal Z\)-stable. The distance and ambient conventions are missing from the supplied statement; the formulation above explicitly uses the customary Kadison--Kastler interpretation, subject to source confirmation.
\subsection{Sources}
\begin{itemize}
\item Local UnsolvedMath registry, permanent ID \(30002218\), source problem code \texttt{catalogue entry}, title: Perturbation Stability of Jiang--Su Absorption for \(C^*\)-Algebras. The supplied record attributes the question to the Oberwolfach Reports Open Problems collection, gives proposed year 2012, and names no individual proposer.
\item Supplied report citation: \emph{\(C^*\)-Algebras, Dynamics, and Classification} (2013), Oberwolfach Reports, pp.\ 3129--3209, DOI \texttt{10.4171/owr/2012/52}. Original OWR question: \url{https://doi.org/10.4171/owr/2012/52}.
\item Supplied literature source, labeled ``perturbations of nuclear C-star algebras'': \url{https://arxiv.org/abs/0910.4953}.
\item Supplied literature source, labeled ``Cuntz semigroup stability under closeness'': \url{https://arxiv.org/abs/1210.4533}.
\item Literature assessment dated exactly \texttt{2026-08-21}: the supplied curation labels the problem partially solved and reports preservation of embeddings of strongly self-absorbing algebras in near-inclusion results and preservation of Cuntz-semigroup data under complete closeness, while stating that preservation of \(\mathcal Z\)-stability in the stated generality remains unproved. This is supplied dated curation, not an independently verified literature assessment; the cited claims and time-sensitive status require independent verification.
\item The registry labels the statement \texttt{corrected\_verified} and reports a manual semantic, typographic, and LaTeX review. These are supplied registry assertions, not mathematical verification of this draft or confirmation of the missing conventions.
\end{itemize}
\end{document}
